\subsection[Complementary elliptic quotients and exact isogeny Selmer groups]{The complementary elliptic quotients and their exact isogeny Selmer groups}

\begin{statement}{A.150}{Proved}
{Actual Prym maps, a conjugation cocycle, and two complete local descents}
Let \(\overline D\) be the proper genus-three curve of A.139,
\[
 aW^4=T(T-1)(T-\tfrac12),\qquad
 a=573q_0^2,\qquad q_0=\frac{M}{2\cdot191^2},
 \qquad M=7\cdot13\cdot23\cdot139\cdot1119169.
\tag{A.150.1}
\]
Put \(K=\Q(\sqrt2)\), \(\epsilon=1+\sqrt2\), and
\[
                 b=1146\epsilon^2M^2,\qquad
                 E_b:\ Y^2=X^3+bX .
\tag{A.150.2}
\]
There are actual degree-two \(K\)-morphisms from
\(\overline D_K\) to \(E_b\) and its conjugate \(E_b^\sigma\),
where \(\sigma(\sqrt2)=-\sqrt2\).  Together with the degree-two
map \(q:\overline D\to E\) of A.143 they induce an isogeny
\[
 \operatorname{Jac}(\overline D)_K
       \longrightarrow E_K\times E_b\times E_b^\sigma .
\tag{A.150.3}
\]
The conjugate-isomorphism
\[
 \mu:E_b^\sigma\longrightarrow E_b,\qquad
             (X,Y)\longmapsto(\epsilon^2X,\epsilon^3Y)
\tag{A.150.4}
\]
satisfies \(\mu\,\sigma(\mu)=[-1]\).

For a nonzero coefficient \(B\in K\), write
\(\mathcal S_B\subset K^\times/K^{\times2}\) for the classes
\(d\) whose smooth proper quartic
\[
                    C_{B,d}:\ N^2=dU^4+(B/d)V^4
\tag{A.150.5}
\]
has a point over every completion of \(K\).
The proof identifies this group with the Selmer group of the
displayed dual two-isogeny.  For (A.150.2) the two exact groups are
\[
 \begin{aligned}
 \mathcal S_b
   =\big\langle&
       3,\ 191,\ 23,\ 139(99+52\sqrt2),\\
       &7(1059+34\sqrt2),\
         7(1059-34\sqrt2)\big\rangle
                   \simeq(\Z/2)^6,\\
 \mathcal S_{-4b}=\mathcal S_{-b}
   &=\langle1119169,-573\rangle\simeq(\Z/2)^2 .
 \end{aligned}
\tag{A.150.6}
\]
The equality of the last two groups uses an explicit
isomorphism preserving the \(X\)-squareclass.
Consequently
\[
                \operatorname{rank}E_b(K)\in\{0,2,4,6\}.
\tag{A.150.7}
\]
The exact groups in (A.150.6) classify the locally soluble
two-isogeny coverings.  The proof does not identify their
classes with global points or assert rank zero.
\end{statement}
\begin{proof}
\emph{Both quotient maps and their inverse field equations.}
For \(\lambda=\pm1/\sqrt2\), the formula
\[
       \iota_\lambda(T,W)
          =\left(\frac{\lambda^2}{T},
                         \frac{\lambda W}{T}\right)
\tag{A.150.8}
\]
preserves (A.150.1) and has order two.  Its invariant functions
\[
 x_\lambda=\frac{aW^2}{T},\qquad
 y_\lambda=\frac{aW(T+\lambda)}T
\tag{A.150.9}
\]
satisfy
\[
       y_\lambda^2=x_\lambda^3+
                         a(\tfrac32+2\lambda)x_\lambda .
\tag{A.150.10}
\]
After cancellation, the equation is precisely
\[
 (T+\lambda)^2=(T-1)(T-\tfrac12)
                                      +(\tfrac32+2\lambda)T .
\]
Conversely, with \(u=x_\lambda^2/a+3/2\),
\[
       T^2-uT+\lambda^2=0,\qquad
       W=\frac{y_\lambda T}{a(T+\lambda)}.
\tag{A.150.11}
\]
Thus the invariant-field extension has degree at most two;
the nonidentity automorphism (A.150.8) fixes both functions,
so its degree is exactly two.  These rational maps extend to
the smooth proper curves.  The target is nonsingular because
its coefficient in (A.150.10) is nonzero.

At \(T=0\), the orders of \(T,W\) are \(4,1\), so
\(x_\lambda,y_\lambda\) have pole orders two and three.
At infinity their orders follow from
\(\operatorname{ord}(T)=-4,\operatorname{ord}(W)=-3\)
and are again minus two and minus three.
Thus these two boundary points map to \(O\).
The points \(T=1,1/2\) map to \((0,0)\).
The involution exchanges each pair, which proves that all
four boundary points are unramified for these quotient maps.

Differentiating the original quartic gives the actual formula
\[
       q_\lambda^*\frac{dx_\lambda}{2y_\lambda}
               =\frac{(T-\lambda)\,dT}{4aW^3}.
\tag{A.150.12}
\]
The degree-four \(T\)-map has four fully ramified branch
points, so Riemann--Hurwitz gives genus three.
The differentials
\[
           \frac{dT}{W^2},\qquad
           \frac{dT}{W^3},\qquad
           \frac{T\,dT}{W^3}
\tag{A.150.13}
\]
are regular: at each finite branch point \(dT\) has order
three in the parameter \(W\), and at infinity their orders
are respectively \(1,4,0\).
They are independent.  The old map in A.143 is
\[
       x=1146(2T-1),\qquad y=8\cdot573^2q_0W^2;
\]
its pulled-back invariant differential is a nonzero multiple
of the first differential in (A.150.13).
The two differentials in (A.150.12) span the other two.
Consequently the map on degree-zero divisor classes given by
the three pushforwards has invertible tangent map and finite
kernel, proving the isogeny (A.150.3).
Restricting its last two coordinates to
\(\ker(q_*)^0\) gives the corresponding isogeny from the
dimension-two Prym factor.  These are the actual divisor
pushforwards, with origins accounted for by degree zero.

For \(\lambda=1/\sqrt2\), scale
\[
           X=(2\cdot191)^2x_\lambda,\qquad
           Y=(2\cdot191)^3y_\lambda.
\tag{A.150.14}
\]
Using \(a=573M^2/(4\cdot191^4)\) gives the coefficient
\(b=1146(1+\sqrt2)^2M^2\) exactly.
The other sign gives \(b^\sigma\).
Since \(b^\sigma=b\epsilon^{-4}\), (A.150.4) follows by
substitution.  Finally \(\epsilon\sigma(\epsilon)=-1\),
so the composite multiplies \(X\) by one and \(Y\) by minus
one.  In particular the actual semilinear group map
\[
 J:E_b(K)\longrightarrow E_b(K),\qquad
 J(X,Y)=(\epsilon^2\sigma X,\epsilon^3\sigma Y)
 \quad\hbox{satisfies}\quad J^2=[-1].
\tag{A.150.15}
\]
This also records the sign of the conjugation cocycle:
every origin-preserving \(K\)-automorphism of this short
Weierstrass model is \((X,Y)\mapsto(X,\pm Y)\), because
the fourth roots of unity in the real field \(K\) are
\(\{\pm1\}\).  Changing \(\mu\) by either automorphism
leaves \(\mu\sigma(\mu)=[-1]\).

\emph{The exact two-isogeny and covering interpretation.}
For \(B\ne0\), put \(E_B:y^2=x^3+Bx\) and
\(E'_B:y^2=x^3-4Bx\).  The two maps are
\[
 \begin{aligned}
 \phi_B:E_B\longrightarrow E'_B,\qquad
     (x,y)&\longmapsto
            \left(x+\frac Bx,\ y(1-B/x^2)\right),\\
 \psi_B:E'_B\longrightarrow E_B,\qquad
     (u,v)&\longmapsto
            \left(\frac{u-4B/u}{4},\
                         \frac{v(1+4B/u^2)}8\right).
 \end{aligned}
\tag{A.150.16}
\]
They extend over their displayed poles, take origins to
origins, and have kernel \(\{O,(0,0)\}\).
Their equations and both compositions \([2]\) follow by
substitution or the tangent duplication formula.
The origin-preserving curve maps are group homomorphisms
by pushforward of the divisor class \([P-O]\).

Define
\[
 \delta_B(O)=1,\quad \delta_B((0,0))=B,\quad
                   \delta_B((x,y))=[x]\ (x\ne0).
\tag{A.150.17}
\]
A line \(y=mx+c\) meets the cubic in three \(x\)-coordinates
whose product is \(c^2\).  This proves multiplicativity
of \(\delta_B\), with tangents obtained by multiplicity.
The formula \(x(P+(0,0))=B/x(P)\) includes the cases
with a zero coordinate.
Its kernel is exactly \(\psi_B E'_B(K)\).
Indeed an image in (A.150.16) has
\(x=(v/(2u))^2\).  Conversely, for \(x=h^2\ne0\) the
literal preimage is
\[
                 u=2x+2y/h,\qquad v=2hu .
\tag{A.150.18}
\]
The equations in (A.150.16) verify its image.
The point \((0,0)\) has a preimage precisely when \(B\)
is a square, in agreement with (A.150.17).

For any \(d\), the quartic (A.150.5) maps to \(E_B\) by
\[
           x=\frac{dU^2}{V^2},\qquad
           y=\frac{dUN}{V^3}.
\tag{A.150.19}
\]
The formulas at \(U=0,V=0\) extend properly.
Over \(K(\sqrt d)\), put \(h=\sqrt d\,U/V\) in
(A.150.18).  This identifies (A.150.19) with \(\psi_B\).
Changing the sign of \(\sqrt d\) changes \(u\) to
\(-4B/u\) and \(v\) to \(4Bv/u^2\), exactly translation
by \((0,0)\) on \(E'_B\).
Thus the quartic is the actual twist by the Kummer
character of \(d\) of this two-isogeny covering.
Its local solubility is equivalent to membership of
\(d\) in the local image of \(\delta_B\), at every place.
Under the identification
\(H^1(K,E'_B[\psi_B])=K^\times/K^{\times2}\), this proves
\[
                   \mathcal S_B
                     =\operatorname{Sel}^{\psi_B}(E'_B/K).
\tag{A.150.20}
\]
In the other direction the isogeny Selmer group for
\(\phi_B\) is \(\mathcal S_{-4B}\).
The scaling
\[
 E_{-b}\longrightarrow E_{-4b},\qquad
                (x,y)\longmapsto(2x,2\sqrt2\,y)
\tag{A.150.21}
\]
preserves the \(x\)-squareclass because \(2=(\sqrt2)^2\).
It proves the stated equality \(\mathcal S_{-4b}=\mathcal S_{-b}\).

\emph{All global candidates and all odd local conditions.}
The ring \(\Z[\sqrt2]\) is norm Euclidean: subtracting
nearest integer coefficients leaves \(r+s\sqrt2\) with
\(|r|,|s|\le1/2\), and
\(|r^2-2s^2|\le1/2<1\).
Its ideals are principal.
Its units are \(\{\pm\epsilon^m:m\in\Z\}\).
For completeness, a positive real unit strictly between
one and \(\epsilon\) would have second integral coefficient
zero or one, from its two real embeddings and norm
\(\pm1\).  The two corresponding norm equations give
no such unit.  Multiplication by a power of \(\epsilon\)
then proves the assertion.
The norm-minus-one unit \(\epsilon\) also makes the
narrow class group trivial.

The complete list of supporting prime generators, all
totally positive, is
\[
 \begin{array}{c|c|c}
 \text{rational prime}&\text{generator or conjugate pair}
                                      &v_{\mathfrak p}(b)\\ \hline
 2&2+\sqrt2&2\\
 3&3&1\\
 191&17\pm7\sqrt2&1\\
 7&3\pm\sqrt2&2\\
 13&13&2\\
 23&5\pm\sqrt2&2\\
 139&139&2\\
 1119169&1059\pm34\sqrt2&2 .
 \end{array}
\tag{A.150.22}
\]
The inert rows \(3,13,139\) have residue fields of size
\(p^2\).  The other odd rows split.
The factorization is the literal equality
\[
 b=(2+\sqrt2)^2\,3\,(17+7\sqrt2)(17-7\sqrt2)\,M^2.
\tag{A.150.23}
\]
At a prime outside this list, \(x\) on \(E_{\pm b}\)
has even valuation: for positive valuation the equation
has leading valuation \(v(x)\), and for negative valuation
it has leading valuation \(3v(x)\).
The covering interpretation therefore restricts every
Selmer class to these twelve prime generators and the
unit classes \(1,-1,\epsilon,-\epsilon\).
Both real embeddings of \(b\) are positive, so the
first curve has only positive \(x\)-classes.
Exactly one of those four unit classes is totally positive.
This gives \(2^{12}=4096\) candidates for \(\mathcal S_b\).
For \(\mathcal S_{-b}\) all four unit classes are retained,
giving \(2^{14}=16384\) candidates.
They all meet the real local condition: one of the two
coefficients \(d,-b/d\) is positive at each real place.

At an odd supporting prime, write
\[
                B=\pi^e B_0,\qquad
                d=\pi^\delta d_0,\qquad
                e\in\{1,2\},\quad\delta\in\{0,1\},
\]
with unit \(B_0,d_0\), for \(B=b\) or \(B=-b\).
Primitive integral \(U,V\) in (A.150.5) give exactly
the following residue-field conditions:
\[
 \begin{array}{c|c|l}
 e&\delta&\text{necessary and sufficient condition}\\ \hline
 1&0&d_0\text{ is a square}\\
 1&1&B_0/d_0\text{ is a square}\\
 2&0&d_0\text{ or }B_0/d_0\text{ is a square}\\
 2&1&-B_0/d_0^2\text{ is a fourth power}.
 \end{array}
\tag{A.150.24}
\]
For \(e=1\), the unequal term valuations force the
unit coefficient in the table to be a square.
For \(e=2,\delta=0\), the two term valuations differ
by two modulo four, so the leading unit is one of
the two listed units.
For \(e=2,\delta=1\), both coordinates must be units;
their valuation-one terms must cancel, giving the
fourth-power condition.
These conditions suffice: the square cases set one
coordinate to zero and lift a square root, and the
last case lifts a nonzero fourth root to a point
with \(N=0\).  Two and four are invertible in these
residue fields, so Hensel's lemma applies.

At any remaining odd prime, the primitive weighted
quartic has smooth genus-one reduction: its quartic
branch polynomial has four distinct roots in the
algebraic closure, since both coefficients are units.
The curve bound
\(q+1-2\sqrt q>0\), from the trace formula and
the Frobenius bound
\cite[(1.5.1), Theorem (1.6)]{AgentDeligneWeilI},
gives a residue-field point; smooth lifting gives a
local point.  Thus no other finite odd conditions
are omitted.

\emph{The exact dyadic calculation and its lifting bound.}
Write \(\varpi=\sqrt2\) in \(K_2=\Q_2(\sqrt2)\).
A unit modulo \(\varpi^5\) is represented by
\(a+c\varpi\), with \(a\bmod8,\ c\bmod4\) and \(a\) odd.
The square-unit residues are exactly
\((1,0)\) and \((3,2)\).
Every element of \(1+\varpi^5\mathcal O_{K_2}\) is a
square, because the polynomial \(Z^2-u\) at one has
error depth at least five, greater than twice the
derivative depth \(v_\varpi(2)=2\).
Consequently the eight unit squareclasses have representatives
\[
        1,\ 1+\varpi,\ 1+2\varpi,\ 1+3\varpi,\
        3+\varpi,\ 3+3\varpi,\ 5,\ 5+2\varpi .
\tag{A.150.25}
\]
Multiplication by \(\varpi^\delta\), \(\delta=0,1\),
gives all sixteen local squareclasses.

The primitive charts \(U=1\), or \(V=1\) with \(U\)
nonunit, exhaust (A.150.5).
In the finite ring
\[
 \mathcal O_{K_2}/\varpi^L
   =\{a+c\varpi:
       a\bmod2^{\lceil L/2\rceil},\
       c\bmod2^{\lfloor L/2\rfloor}\},
\tag{A.150.26}
\]
multiplication is
\((a,c)(a',c')=(aa'+2cc',ac'+a'c)\).
Exhausting these charts excludes eight classes for
\(b\) already at \(L=5\); for \(-b\), \(L=7\)
leaves only the two classes below.
At \(L=13\) the complete surviving sets are
\[
 \begin{aligned}
 b:&\quad
 \{\varpi^\delta u:\ \delta=0,1,\
                 u\in\{1,1+2\varpi,5,5+2\varpi\}\},\\
 -b:&\quad \{1,1+2\varpi\}.
 \end{aligned}
\tag{A.150.27}
\]
The exhaustive finite ring loops are retained in the
certificate; the following are actual residue witnesses,
written as \((U,V,N)\) modulo \(\varpi^{13}\):
\[
 \begin{array}{c|c|c}
 B&d&(U,V,N)\\ \hline
 b&1&(1,0,1)\\
 b&1+2\varpi&(\varpi,1,16+13\varpi)\\
 b&5&(0,1,10+21\varpi)\\
 b&5+2\varpi&(1,1,21+\varpi)\\
 b&\varpi&(1,1+3\varpi,10+6\varpi)\\
 b&\varpi(1+2\varpi)&(1,1+2\varpi,14+2\varpi)\\
 b&5\varpi&(1,1,14+4\varpi)\\
 b&\varpi(5+2\varpi)&(1,1+3\varpi,10+4\varpi)\\
 -b&1&(1,0,1)\\
 -b&1+2\varpi&(1,1,31+11\varpi).
 \end{array}
\tag{A.150.28}
\]
These are sufficient, rather than merely finite
congruence tests.  Here \(v_\varpi(b)=2\).
For \(F=N^2-dU^4-(B/d)V^4\), a unit \(U\)
gives
\(v_\varpi(F_U)=4+\delta\le5\).
When \(U\) is nonunit, \(V\) is a unit and
\(v_\varpi(F_V)=6-\delta\le6\).
In either case the error depth thirteen is
strictly greater than twice the derivative depth.
One-variable Hensel lifting preserves primitiveness
and gives an actual local point.
This proves the exact local sets (A.150.27).

\emph{The finite global intersection and the rank implication.}
Enumerate the twelve prime bits in the order of
(A.150.22), taking the positive sign of each
conjugate pair first.
At a split prime generated by \(A+C\sqrt2\),
use the actual residue map
\(\sqrt2\mapsto-A/C\bmod p\).
At an inert prime use
\(\F_p[t]/(t^2-2)\).
Omit a prime's own generator to compute \(B_0,d_0\)
in (A.150.24); all other factors are nonzero.
Testing \(u^{(q-1)/2}=1\) for squares and
\(u^{(q-1)/\gcd(4,q-1)}=1\) for fourth powers
then gives a finite, complete calculation.
Its outcome is
\[
 \begin{array}{c|r|r|r}
 &\text{all global candidates}&\text{odd local conditions}
                                      &\text{all local conditions}\\ \hline
 b&4096&64&64\\
 -b&16384&8&4 .
 \end{array}
\tag{A.150.29}
\]
The complete resulting bit spans, including the two
unit bits in the second row, have bases
\[
 \begin{array}{c|l}
 b&2,\ 12,\ 384,\ 644,\ 1072,\ 2096\\
 -b&3072,\ 4110 .
 \end{array}
\tag{A.150.30}
\]
Here the minus-one bit is \(4096\) and the
\(\epsilon\) bit is \(8192\).
Literal multiplication of the generators converts
(A.150.30) to exactly (A.150.6).
The certificate verifies equality of the entire
spans with the enumerated survivors, including
their independence.  All tests are now necessary
and sufficient at every place, so this computes
the Selmer groups themselves.

Both \(b\) and \(-b\) are nonsquares in \(K\),
as their valuations at three are odd.
Thus both elliptic curves have exactly two
rational two-torsion points.
The nonzero kernel point of \(\psi_b\) is not in
the image of \(\phi_b\), since its descent class
is \(-4b\), a nonsquare.
Applying the two kernels and their quotient maps
in (A.150.16)--(A.150.18) gives
\[
 2^{\operatorname{rank}E_b(K)}
       =\frac{|\operatorname{im}\delta_b|\,
                   |\operatorname{im}\delta_{-4b}|}{4}
       \le\frac{64\cdot4}{4}=64 .
\tag{A.150.31}
\]
The finite-generation input is the number-field
Mordell--Weil theorem
\cite[Chapter IV, unnumbered finite basis theorem,
printed p.~105, and \S4]{AgentMilneEC2021}.
Finally (A.150.15) makes
\(E_b(K)\otimes_\Z\Q\) a vector space over
\(\Q[J]\simeq\Q(i)\).  Its \(\Q\)-dimension is even,
proving (A.150.7).
\end{proof}

The exact programs retaining the morphism identities,
every residue witness, and every global squareclass are
\begin{flushleft}
\path{code/harmonic_prym_quotient_certificate.py}\\
\path{code/harmonic_prym_dyadic_filter.py}\\
\path{code/harmonic_prym_odd_filter.py}.
\end{flushleft}
The consolidated
\path{code/harmonic_prym_selmer_certificate.py}
reruns them and independently verifies the literal
Hensel inequalities, both isogeny maps, their composition,
the conjugation cocycle, and the complete spans.
Its output and data are retained beside the source.
